% !TeX root = ../../Book2.tex
% 纲要标题：### 10.3 Witt 理论初步
% 纲要位置：计划纲要.md，第 1386 行。

\section{Witt 理论初步}
\label{b2:sec:1003}

从本节起固定 \(\operatorname{char}k\ne2\)，所讨论的二次空间均有限维，且以 \(b=\beta_q/2\) 为相应对称型。实数域上的正负号分类不能直接搬到一般域上，但非零零点仍能帮助我们分离出一种固定的两维型。Witt 理论先取出这些两维部分，再研究剩下没有非零零点的部分。

% 纲要要点（供目录核对）：
%
% * 双曲平面、各向同性与各向异性；各向同性不等于属于根空间
% * Witt 分解与消去定理；反射用于对齐共同非退化分量
% * Witt 群的概念；忽略双曲部分后由相反型给出加法逆元
%

\subsection{双曲平面、各向同性与各向异性}
\label{b2:subsec:1003:01}

\begin{definition}\label{b2:def:isotropic-hyperbolic}
设 \(k\) 为特征不为二的域，\((V,q)\) 为 \(k\) 上的二次空间。若 \(v\in V\) 满足 \(v\ne0\) 且 \(q(v)=0\)，就称 \(v\) 为\term[gexiang-tongxing-xiangliang]{各向同性向量}[isotropic vector]。若不存在这样的向量，就称 \((V,q)\)\term[gexiang-yixing]{各向异性}[anisotropic]；零空间按此定义也是各向异性的。若子空间 \(T\subseteq V\) 满足 \(q|_T=0\)，就称 \(T\) 为\term[quan-gexiang-tongxing]{全各向同性子空间}[totally isotropic subspace]。
\end{definition}

各向同性只要求一个非零向量的自身取值 \(q(v)=b(v,v)\) 为零；属于根空间则要求 \(b(v,w)=0\) 对所有 \(w\in V\) 成立。因此根空间中的非零向量必各向同性，反过来却不成立。全各向同性要求子空间中每个向量的二次取值为零，由极化公式，这等价于 \(b|_{T\times T}=0\)，比仅指定若干基向量各自取值为零更强。非退化空间的根空间为零，但仍可以有非零全各向同性子空间，例如上一节的 \(x^2-y^2\)。

\begin{definition}\label{b2:def:hyperbolic-plane}
设 \(k\) 为特征不为二的域，\((V,q)\) 为二维 \(k\) 二次空间，\(b=\beta_q/2\)。如果 \(V\) 有基 \(e,f\)，满足 \(q(e)=q(f)=0\)、\(b(e,f)=1\)，就称 \((V,q)\) 为\term[shuangqu-pingmian]{双曲平面}[hyperbolic plane]，记为 \(H\)。它的表达式及矩阵为
\[
q(xe+yf)=2xy,\qquad B_H=\begin{pmatrix}0&1\\1&0\end{pmatrix}.
\]
若一个 \(k\) 二次空间等距于若干个 \(H\) 的正交直和，就称其为\term[shuangqu-kongjian]{双曲空间}[hyperbolic space]，并记 \(H^{\perp r}\) 为 \(r\) 个双曲平面的正交直和。
\end{definition}
\symidx{\(H\)、\(H^{\perp r}\)：双曲平面及其正交直和}

表达式中的因子二来自本书的约定 \(q(v)=b(v,v)\)：两个混合项分别为 \(xyb(e,f)\) 与 \(xyb(f,e)\)，各等于 \(xy\)。若把未归一化的极化型 \(\beta_q\) 在基向量间的值取为一，才会得到 \(xy\) 的表达式。

\ProseParagraphBreak
所有这样给出的平面都等距，因为把指定的 \(e,f\) 相互对应便保持矩阵。\(\det B_H=-1\ne0\)，所以它非退化，却包含全各向同性直线 \(ke\) 与 \(kf\)。这两条直线的和不是全各向同性子空间，因为 \(q(e+f)=2\ne0\)；各向同性向量 \(e\) 也不属于根空间，因为 \(b(e,f)=1\)。基 \(e+f/2,e-f/2\) 相互正交，二次取值为一和负一，因而
\[
H\cong\langle1,-1\rangle.
\]
若 \(k=\mathbb R\)，这正是一个正方向与一个负方向组成的部分。

\begin{lemma}[双曲平面的分离]\label{b2:lem:hyperbolic-plane-split}
设 \(k\) 为特征不为二的域，\((V,q)\) 为有限维非退化 \(k\) 二次空间，\(b=\beta_q/2\)。若 \(e\in V\) 满足 \(e\ne0\) 且 \(q(e)=0\)，则存在 \(f\in V\)，使 \(e,f\) 生成一个双曲平面 \(H\)，并有 \(V=H\perp H^\perp\)，其中正交补非退化。
\end{lemma}
\begin{proof}
由于 \(q(e)=0\)，直线 \(ke\) 上的限制型是零型，不能把它单独取作非退化正交分量，也不能像分离非零对角项时那样除以 \(q(e)\)。我们要给它补上另一个各向同性方向，使两个方向之间的配对非零。

非退化性保证某个 \(f_0\) 满足 \(b(e,f_0)\ne0\)。缩放以后令此值为一。先尝试 \(f=f_0+te\)，其中 \(t\in k\)。由 \(q(v)=b(v,v)\)、\(q(e)=0\) 及 \(b(e,f_0)=b(f_0,e)=1\)，有
\[
b(e,f)=1,\qquad
q(f)=q(f_0)+2t\,b(f_0,e)+t^2q(e)=q(f_0)+2t.
\]
所以保持 \(b(e,f)=1\) 的同时，让 \(q(f)=0\) 只需求解线性方程 \(q(f_0)+2t=0\)。因为特征不为二，取 \(t=-q(f_0)/2\)，即
\[
f=f_0-\frac{q(f_0)}2e.
\]
两向量的配对矩阵为 \(B_H\)，所以它们线性无关且限制型非退化。现在应用\cref{b2:thm:orthogonal-complement}，即得所需正交分解与正交补的非退化性。
\end{proof}

\begin{proposition}\label{b2:prop:isotropic-subspace-hyperbolic}
设 \(k\) 为特征不为二的域，\((V,q)\) 为有限维非退化 \(k\) 二次空间。若 \(T\subseteq V\) 为 \(d\) 维全各向同性子空间，则 \(T\) 包含于某个等距于 \(H^{\perp d}\) 的非退化子空间。
\end{proposition}
\begin{proof}
对 \(d\) 归纳，\(d=0\) 时取零空间。\(d>0\) 时，取 \(0\ne e\in T\)，由\cref{b2:lem:hyperbolic-plane-split}补上 \(f\)，令 \(H=ke+kf\)。因为 \(T\) 全各向同性，其中每个 \(t\) 已与 \(e\) 正交；要让它落入 \(H^\perp\)，还须消去它与 \(f\) 的配对。利用 \(b(e,f)=1\)，沿 \(e\) 调整，置
\[
t'=t-b(t,f)e.
\]
此时 \(b(t',e)=0\)，且 \(b(t',f)=b(t,f)-b(t,f)b(e,f)=0\)。又 \(t'\in T\)，所以每个 \(t\) 都有表示 \(t=b(t,f)e+t'\)。由于 \(e\notin H^\perp\)，这个表示给出直和
\[
T=ke\oplus\bigl(T\cap H^\perp\bigr).
\]
交的维数为 \(d-1\)，仍全各向同性；在非退化的 \(H^\perp\) 内应用归纳假设，再加上 \(H\)，便得到所求子空间。
\end{proof}

\subsection{Witt 分解与消去定理}
\label{b2:subsec:1003:02}

不断分离双曲平面，每次维数减少二，便能得到分解的存在性。要说明所得各向异性部分不依赖选择，需要先证明正交直和可以消去共同的非退化部分。这里的困难是，\(U\perp V\) 到 \(U\perp W\) 的等距同构未必把指定的 \(U\) 分量送到另一边的 \(U\) 分量。将 \(U\) 对角化后，只需逐条对齐共同的非退化直线，再把映射限制到正交补；下面的反射正是用来调整这些直线的像。

\begin{lemma}[反射与等长向量]\label{b2:lem:orthogonal-reflection}
设 \(k\) 为特征不为二的域，\((V,q)\) 为有限维 \(k\) 二次空间，\(b=\beta_q/2\)。若 \(z\in V\) 满足 \(q(z)\ne0\)，则 \(V=kz\perp(kz)^\perp\)。在 \(kz\) 上取负恒等、在 \((kz)^\perp\) 上取恒等所定义的映射 \(\tau_z\) 是等距自同构，称为沿 \(z\) 的\term[zhengjiao-fanshe]{反射}[reflection]。它的公式是
\[
\tau_z(v)=v-\frac{2b(v,z)}{q(z)}z.
\]
若 \(u,v\in V\) 满足 \(q(u)=q(v)=a\ne0\)，则存在至多两个这样的反射的复合，将 \(u\) 送到 \(v\)。
\end{lemma}
\begin{proof}
因为 \(q(z)\ne0\)，直线 \(kz\) 的限制型非退化，\cref{b2:thm:orthogonal-complement}给出 \(V=kz\perp(kz)^\perp\)，这里不要求整个 \(V\) 非退化。对任意 \(v\in V\)，这一分解具体写成
\[
v=\frac{b(v,z)}{q(z)}z+v_\perp,
\qquad v_\perp\in(kz)^\perp.
\]
只把第一项取负，得到 \(-\dfrac{b(v,z)}{q(z)}z+v_\perp=v-\dfrac{2b(v,z)}{q(z)}z\)，这就导出了所给公式。两个正交分量上的作用均保持二次取值，且平方为恒等，所以 \(\tau_z\) 是等距自同构。

对等长向量，计算得
\[
q(u-v)=2\bigl(a-b(u,v)\bigr),\qquad
q(u+v)=2\bigl(a+b(u,v)\bigr).
\]
两者之和为 \(4a\ne0\)，故至少一个非零。若 \(q(u-v)\ne0\)，则
\[
\frac{2b(u,u-v)}{q(u-v)}=1,
\qquad \tau_{u-v}(u)=v.
\]
若 \(q(u-v)=0\)，则 \(q(u+v)\ne0\)，同样计算得 \(\tau_{u+v}(u)=-v\)。再施加 \(\tau_v\)，便把 \(-v\) 送到 \(v\)。这里 \(q(v)=a\ne0\)，所以第二个反射确有定义。\(u=v\) 时也可直接取恒等映射。
\end{proof}

这项反射构造可与 Roman\cite[Chapter 11, Theorem 11.32]{Roman2008AdvancedLinearAlgebra}对照。非零长度的条件十分具体：它既保证至少一个 \(u\pm v\) 能定义反射，也保证必要时能再沿 \(v\) 反射。各向同性向量不能直接使用这个论证。

\begin{theorem}[Witt 消去]\label{b2:thm:witt-cancellation}
设 \(k\) 为特征不为二的域，\(U,V,W\) 为有限维非退化 \(k\) 二次空间。若
\[
U\perp V\cong U\perp W,
\]
则 \(V\cong W\)。
\end{theorem}
\begin{proof}
先设共同部分一维。在两边分别取基向量 \(e,f\)，使 \(q(e)=q(f)=a\ne0\)，并给定等距同构
\[
\phi:ke\perp V\longrightarrow kf\perp W.
\]
目标空间中的 \(\phi(e)\) 与 \(f\) 等长且长度非零。由\cref{b2:lem:orthogonal-reflection}，有目标空间的等距自同构 \(\rho\)，使 \(\rho\phi(e)=f\)。因此 \(\rho\phi\) 将 \((ke)^\perp=V\) 送入 \((kf)^\perp=W\)：对 \(v\in V\)，保持配对给出 \(b\bigl(\rho\phi(v),f\bigr)=b(v,e)=0\)。反向应用它的逆同构，得到像恰为 \(W\)，故限制映射是 \(V\to W\) 的等距同构。

一般情形中，由\cref{b2:thm:quadratic-diagonalization}，共同非退化部分可写成
\[
U=\langle a_1\rangle\perp\cdots\perp\langle a_d\rangle,
\qquad a_i\ne0.
\]
依次对两边共同的一维部分应用已证情形。每次余下的空间仍是非退化正交直和，故下一次消去仍满足假设。经过 \(d\) 次即得 \(V\cong W\)；\(d=0\) 时本来就是该结论。
\end{proof}

上述结果称为\term[Witt-xiaoqu-dingli]{Witt 消去定理}[Witt cancellation theorem]。

\begin{theorem}[Witt 分解]\label{b2:thm:witt-decomposition}
设 \(k\) 为特征不为二的域，\((V,q)\) 为有限维非退化 \(k\) 二次空间，则 \(V\) 有分解
\[
V\cong H^{\perp r}\perp A,
\]
其中 \(A\) 各向异性。整数 \(r\) 与 \(A\) 的等距类型唯一。\(r\) 等于全各向同性子空间的最大维数，称为\term[Witt-zhishu]{Witt 指数}[Witt index]。
\end{theorem}
\begin{proof}
若有非零各向同性向量，便用\cref{b2:lem:hyperbolic-plane-split}取出一个双曲平面。在非退化正交补上重复，每次维数减少二，因此有限步后剩余部分 \(A\) 没有非零各向同性向量。若原空间已经各向异性，则取 \(r=0\)。

设还有 \(V\cong H^{\perp s}\perp B\)，其中 \(B\) 各向异性。不妨 \(r\ge s\)。由\cref{b2:thm:witt-cancellation}消去 \(s\) 个双曲平面，得到 \(H^{\perp(r-s)}\perp A\cong B\)。若 \(r>s\)，左侧含非零各向同性向量，右侧却没有，矛盾。因此 \(r=s\)，继续用所得等距即得 \(A\cong B\)。

所列双曲平面中的全部 \(e_i\) 线性无关，且 \(q(e_i)=0\)、不同平面的向量之间配对为零，所以它们生成一个 \(r\) 维全各向同性子空间。反过来，任意 \(d\) 维全各向同性子空间由\cref{b2:prop:isotropic-subspace-hyperbolic}包含在某个 \(H^{\perp d}\) 中。将其非退化正交补再作已经证明存在的 Witt 分解，得到整个 \(V\) 的一个分解，其中双曲平面的数目至少为 \(d\)。由唯一性，它恰为 \(r\)，故 \(d\le r\)。
\end{proof}

这一分解称为\term[Witt-fenjie]{Witt 分解}[Witt decomposition]。若 \(q\) 退化，记 \(R=\operatorname{rad}b\)。对 \(z\in R\)，有 \(q(z)=b(z,z)=0\) 及 \(b(v,z)=0\)，所以 \(q(v+z)=q(v)\)。因此商空间上有良定义的二次型
\[
\bar q(v+R)=q(v).
\]
它非退化：若 \(v+R\) 与全部商向量正交，就有 \(b(v,w)=0\) 对所有 \(w\in V\) 成立，从而 \(v\in R\)、\(v+R=0\)。根空间 \(R\)、商空间 \(V/R\) 及这个商型都由原型直接确定，不需要选取补空间。

\ProseParagraphBreak
要在 \(V\) 内写出直和分解，可取 \(R\) 的一个普通线性补空间 \(V_0\)。自然映射 \(V_0\to V/R\) 是线性同构，且保持二次取值，所以 \(V_0\) 的限制型非退化。对它应用 Witt 分解，得到
\[
V=R\perp V_0
\cong R\perp H^{\perp r}\perp A.
\]
第一项是零型。与第一章的商空间和补空间之别相同，\(V_0\) 在 \(V\) 中的具体位置一般不唯一，但每个这样的补空间都等距于同一个商型；商空间的 Witt 指数与各向异性部分的等距类型因而由原二次型确定。

\subsection{Witt 群的概念}
\label{b2:subsec:1003:03}

正交直和给出的等距类可以相加，但对非零 \(V\)，\((V,q)\perp(V,-q)\) 的底层空间是 \(V\oplus V\)，维数仍为 \(2\dim V\)，不能等距于零空间。一维情形已经揭示了应忽略什么：对 \(a\in k^\times\)，在 \(\langle a,-a\rangle\) 中取
\[
e=(1,1),\qquad f=\left(\frac1{2a},-\frac1{2a}\right).
\]
直接计算得 \(q(e)=q(f)=0\)、\(b(e,f)=1\)。这两个向量线性无关，因为把 \(ce+df=0\) 分别与 \(e,f\) 配对便得 \(d=c=0\)。所以 \(\langle a,-a\rangle\cong H\)。若允许增删双曲平面而不改变类，这对互为负号的一维型之和便为零；对角化后同样的配对将使 \((V,-q)\) 成为 \((V,q)\) 的加法逆元。

\begin{definition}\label{b2:def:witt-equivalence}
设 \(k\) 为特征不为二的域，\(V,W\) 为有限维非退化 \(k\) 二次空间。如果存在整数 \(r,s\ge0\)，使
\[
V\perp H^{\perp r}\cong W\perp H^{\perp s},
\]
就称 \(V\) 与 \(W\)\term[Witt-dengjia]{Witt 等价}[Witt equivalent]。所有等价类组成的集合记为 \(W(k)\)。在类上定义 \([V]+[W]=[V\perp W]\)。
\end{definition}

\begin{theorem}\label{b2:thm:witt-group}
设 \(k\) 为特征不为二的域，\(W(k)\) 为有限维非退化 \(k\) 二次空间按添加双曲平面所得的 Witt 等价类集合，则 Witt 等价是等价关系，正交直和给出 \(W(k)\) 上良定义的交换群运算。零元为零空间的类，\((V,q)\) 的负元由 \((V,-q)\) 表示。每个类恰有一个各向异性代表的等距类型。
\end{theorem}
\begin{proof}
自反性和对称性直接成立。若 \(V\perp H^{\perp r}\cong W\perp H^{\perp s}\)，又 \(W\perp H^{\perp u}\cong X\perp H^{\perp v}\)，则给第一个等距加上 \(H^{\perp u}\)，给第二个加上 \(H^{\perp s}\)，得到
\[
V\perp H^{\perp(r+u)}\cong X\perp H^{\perp(v+s)}.
\]
故关系传递。对两个等距分别取正交直和，便证明加法与代表选择无关。直和的结合、交换及零空间的单位性质给出相应的群公理。

还须构造负元。将 \(q\) 对角化为 \(\langle a_1,\ldots,a_n\rangle\)，其中 \(a_i\ne0\)。\(q\perp(-q)\) 重排分量后为若干个 \(\langle a_i,-a_i\rangle\)，由定义前的一维计算，每个分量都等距于 \(H\)。因此 \((V,q)\perp(V,-q)\cong H^{\perp n}\)，其类为零，给出负元。

\cref{b2:thm:witt-decomposition}说明每个类都由各向异性部分表示。若各向异性 \(A,B\) Witt 等价，则 \(A\perp H^{\perp r}\cong B\perp H^{\perp s}\) 的两边已是 Witt 分解，唯一性迫使 \(r=s\) 且 \(A\cong B\)。
\end{proof}

这个群称为域 \(k\) 的\term[Witt-qun]{Witt 群}[Witt group]。记号中的集合没有大小困难：每个有限维型都可以由某个有限对称矩阵表示，全部这些矩阵组成集合，再取等价类即可。
\symidx{\(W(k)\)：特征不为二的域的 Witt 群}

\begin{proposition}\label{b2:prop:witt-real-complex}
对有限维非退化实二次空间，正惯性指数与负惯性指数之差给出群同构 \(W(\mathbb R)\cong\mathbb Z\)。对有限维非退化复二次空间，维数的奇偶性给出群同构 \(W(\mathbb C)\cong\mathbb Z/2\mathbb Z\)。
\end{proposition}
\begin{proof}
实型由\cref{b2:thm:sylvester-inertia}分为 \(p_+\) 个 \(\langle1\rangle\) 与 \(p_-\) 个 \(\langle-1\rangle\)。一正一负合成双曲平面；配对后剩余系数全为一或全为负一，是各向异性型。它的维数与符号由差 \(p_+-p_-\) 决定，所以其 Witt 类只由这个差决定。符号差在正交直和下相加，加双曲平面不改变它；差为零时全部分量都能配成双曲平面，而任意整数都能由全正或全负型实现，故所得同态既单射又满射。

复数域上，先由\cref{b2:thm:quadratic-diagonalization}对角化，再利用每个非零复数都有平方根，把全部非零系数缩放为一。因此非退化型只有维数这一分类数据。\(\langle1,1\rangle\cong\langle1,-1\rangle\cong H\)，第一个等距通过一个坐标乘以 \(i\) 得到。因此偶数维型的类为零，奇数维型的类为 \(\bigl[\langle1\rangle\bigr]\)。加双曲平面不改变维数奇偶，故这个非零类不可能等于零，得到所述二阶群。
\end{proof}

张量积还能在这个群上定义乘法。\appref{b2:app:B}的\cref{b2:thm:B-witt-ring}证明该乘法对Witt类良定义，\cref{b2:prop:B-discriminant-witt}进一步讨论维数奇偶、判别式和基本理想；这些是本节加法理论之后的补充。
